%!TEX root = lebenslauf.tex
\begin{quote}
	\begin{tabular}{p{4cm}l}
		01/2007 - 12/2008&Mobile Code (MobCo)\\
		&Adaptive Architekturen f"ur mobile Anwendungen\\\\
		&\begin{minipage}{10cm}
		\begin{tabbing}
			Technologien:\ \=\kill
			Realisierung:\>Prof. f. Angw. Telematik/E-Business\\
			\>Universit"at Leipzig\\ 
			Typ:\>Forschungsprojekt\\
			Technologien:\>Java, Java ME, OSGi, R-OSGI, db4o
			\\Stakeholder:\>Deutsche Telekom Laboratories,\\
			\>Universit"at Leipzig
			\\Rolle:\>Mitarbeiter\\
			Aufgaben:\>Entwicklung
		\end{tabbing}
		In Zusammenarbeit mit den Deutsche Telekom Laboratories wurde mit besonderem Fokus auf mobile Informationssysteme im Projekt MobCo eine Middleware zur Erstellung und den Betrieb von adaptiven mobilen Anwendungen erstellt. Die MobCo eingesetzten Technologien umfassen Java, Java ME, das OSGi-Framework (ein leichtgewichtiges Komponentenframework und Diensteplattform), das darauf aufbauende R-OSGi (mit dem verteilte OSGi-Dienste kommunizieren k"onnen) und db4o (eine objektorientierte Datenbank).
		\end{minipage}
	\end{tabular}
\end{quote}